\documentclass[aspectratio=169,10pt]{beamer}
\makeatletter\def\input@path{{../estilo/}{estilo/}}\makeatother
\usepackage{devmobile}
\modulo{11}{Testes e Qualidade}{Pirâmide de testes, testes unitários com coroutines, fakes e mocks, testes de UI em Compose, desempenho, acessibilidade e integração contínua}

\begin{document}
\capa

\begin{frame}{Objetivos do módulo}
  \begin{conceito}[Ao final deste módulo você será capaz de]
    \begin{itemize}
      \item Planejar testes segundo a pirâmide de testes.
      \item Testar funções, casos de uso e ViewModels com JUnit, coroutines e Turbine.
      \item Usar \emph{fakes} e \emph{mocks} para isolar dependências.
      \item Escrever testes de interface em Compose e testes de banco com Room.
      \item Diagnosticar problemas de desempenho e recomposição.
      \item Tornar o app acessível e automatizar a verificação com integração contínua.
    \end{itemize}
  \end{conceito}
\end{frame}

\sumario

% =============================================================================
\section{Estratégia de testes}

\begin{frame}{A pirâmide de testes}
  \begin{columns}[c]
    \begin{column}{0.5\textwidth}
      \centering
      \begin{tikzpicture}[x=0.9cm]
        \fill[petroleo] (-3,0) -- (3,0) -- (2,1.1) -- (-2,1.1) -- cycle;
        \fill[petroleoesc] (-2,1.15) -- (2,1.15) -- (1,2.25) -- (-1,2.25) -- cycle;
        \fill[ambar] (-1,2.3) -- (1,2.3) -- (0,3.4) -- cycle;
        \node[white,font=\scriptsize\bfseries,align=center] at (0,0.55) {Unitários\\ \tiny JVM, milissegundos};
        \node[white,font=\scriptsize\bfseries,align=center] at (0,1.7) {Integração\\ \tiny Room, repositórios};
        \node[white,font=\tiny\bfseries,align=center] at (0,2.75) {UI / E2E};
        \draw[setaambar] (3.4,0.2) -- (3.4,3.2) node[rotulo,right,pos=0.5,align=left]{mais lentos,\\ caros e\\ frágeis};
      \end{tikzpicture}
    \end{column}
    \begin{column}{0.48\textwidth}
      \small
      \begin{itemize}
        \item \textbf{Unitários} (\arq{src/test}): muitos, rápidos, rodam no computador.
        \item \textbf{Integração}: várias peças juntas (DAO + banco em memória).
        \item \textbf{UI / ponta a ponta} (\arq{src/androidTest}): poucos, verificam fluxos críticos no aparelho.
      \end{itemize}
      \begin{dica}[O que testar primeiro]
        Regras de negócio (casos de uso), transformações de dados e os ViewModels. É onde estão os erros mais caros.
      \end{dica}
      \begin{conceito}[Arquitetura testável]
        As camadas, as interfaces e a injeção de dependência dos Módulos 06--07 existem justamente para permitir isso.
      \end{conceito}
    \end{column}
  \end{columns}
\end{frame}

% =============================================================================
\section{Testes unitários}

\begin{frame}[fragile]{Primeiro teste: funções puras}
  \begin{columns}[T]
    \begin{column}{0.54\textwidth}
\begin{codigo}[title=src/test/.../ConversoesTest.kt]{Kotlin}
class ConversoesTest {

  @Test
  fun `zero celsius é 32 fahrenheit`() {
    assertEquals(32.0, 0.0.cParaF(), 0.001)
  }

  @Test
  fun `adc de 12 bits no fundo de escala vale vref`() {
    assertEquals(3.3, 4095.adcParaVolts(), 0.001)
  }

  @Test
  fun `leitura fora da faixa gera alerta`() {
    val alertas = VerificadorFaixa(10.0..40.0)
      .verificar(listOf(25.0, 41.2, 9.0))
    assertEquals(listOf(41.2, 9.0), alertas)
  }
}
\end{codigo}
    \end{column}
    \begin{column}{0.44\textwidth}
      \begin{conceito}[Estrutura AAA]
        \textbf{Arrange} (prepare os dados), \textbf{Act} (execute) e \textbf{Assert} (verifique). Um comportamento por teste.
      \end{conceito}
      \begin{dica}[Nomes legíveis]
        Kotlin aceita nomes de função entre crases: o relatório de testes vira uma especificação em português.
      \end{dica}
      \begin{terminal}[title=Executando]
./gradlew test
# relatório em app/build/reports/tests/
      \end{terminal}
    \end{column}
  \end{columns}
\end{frame}

\begin{frame}[fragile]{Testando coroutines}
  \begin{columns}[T]
    \begin{column}{0.54\textwidth}
\begin{codigo}[title=MainDispatcherRule.kt]{Kotlin}
// Substitui Dispatchers.Main (inexistente na JVM)
class MainDispatcherRule(
  val dispatcher: TestDispatcher =
    UnconfinedTestDispatcher()
) : TestWatcher() {
  override fun starting(d: Description) =
    Dispatchers.setMain(dispatcher)
  override fun finished(d: Description) =
    Dispatchers.resetMain()
}
\end{codigo}
\begin{codigo}[title=Teste com runTest]{Kotlin}
@Test
fun `sincronizar tenta de novo após falha`() = runTest {
  val api = FakeApi(falhasAntesDeResponder = 2)
  val repo = BancadaRepositoryImpl(FakeDao(), api)
  repo.sincronizarComRetry()   // delays são "pulados"
  assertEquals(3, api.chamadas)
}
\end{codigo}
    \end{column}
    \begin{column}{0.44\textwidth}
      \begin{conceito}[\cod{runTest}]
        Executa código \cod{suspend} em testes com \textbf{tempo virtual}: um \cod{delay(60\_000)} termina instantaneamente.
      \end{conceito}
      \begin{dica}[Dispatchers injetados]
        Repositórios que recebem o dispatcher no construtor (Módulo 03) podem receber um \cod{StandardTestDispatcher} no teste.
      \end{dica}
      \begin{dica}[Dependência]
        \cod{kotlinx-coroutines-test} em \cod{testImplementation}.
      \end{dica}
    \end{column}
  \end{columns}
\end{frame}

\begin{frame}[fragile]{Fakes \emph{versus} mocks}
  \begin{columns}[T]
    \begin{column}{0.5\textwidth}
\begin{codigo}[title=Fake: implementação simples]{Kotlin}
class FakeBancadaRepository : BancadaRepository {
  private val dados =
    MutableStateFlow<List<Bancada>>(emptyList())
  var falharSync = false

  fun emitir(lista: List<Bancada>) { dados.value = lista }

  override fun bancadas() = dados
  override suspend fun sincronizar() {
    if (falharSync) throw IOException("sem rede")
  }
  override suspend fun reservar(id: Int, ate: Instant) {}
}
\end{codigo}
    \end{column}
    \begin{column}{0.48\textwidth}
\begin{codigo}[title=Mock com MockK]{Kotlin}
val api = mockk<LabApi>()
coEvery { api.listarBancadas() } throws IOException()

val repo = BancadaRepositoryImpl(FakeDao(), api)
val r = runCatching { repo.sincronizar() }

assertTrue(r.isFailure)
coVerify(exactly = 1) { api.listarBancadas() }
\end{codigo}
      \begin{conceito}[Quando usar]
        \textbf{Fake}: reutilizável, comportamento realista; preferido pelo Google. \textbf{Mock}: verificar \emph{interações} (foi chamado? quantas vezes?).
      \end{conceito}
    \end{column}
  \end{columns}
\end{frame}

\begin{frame}[fragile]{Testando o ViewModel com Turbine}
  \begin{columns}[T]
    \begin{column}{0.58\textwidth}
\begin{codigo}[title=PainelViewModelTest.kt]{Kotlin}
class PainelViewModelTest {
  @get:Rule val main = MainDispatcherRule()
  private val repo = FakeBancadaRepository()
  private val vm = PainelViewModel(repo)

  @Test
  fun `filtro mostra apenas bancadas correspondentes`() =
    runTest {
      repo.emitir(listOf(Bancada(1, "Eletrônica", 24.0),
                         Bancada(2, "Máquinas", 22.0)))
      vm.uiState.test {
        assertEquals(2, awaitItem().bancadas.size)
        vm.onFiltroChange("Máq")
        val estado = awaitItem()
        assertEquals(1, estado.bancadas.size)
        assertEquals("Máquinas", estado.bancadas[0].titulo)
        cancelAndIgnoreRemainingEvents()
      }
    }
}
\end{codigo}
    \end{column}
    \begin{column}{0.4\textwidth}
      \begin{conceito}[Turbine]
        Biblioteca da Cash App que coleta um \cod{Flow} e permite verificar cada emissão em sequência com \cod{awaitItem()}.
      \end{conceito}
      \begin{dica}
        Teste o \textbf{estado} que o ViewModel expõe, não seus detalhes internos. Assim o teste sobrevive a refatorações.
      \end{dica}
      \begin{atencao}
        \cod{stateIn(WhileSubscribed)} só calcula com alguém coletando. O \cod{test \{\}} do Turbine é esse coletor.
      \end{atencao}
    \end{column}
  \end{columns}
\end{frame}

% =============================================================================
\section{Testes de integração e de UI}

\begin{frame}[fragile]{Room com banco em memória}
  \begin{columns}[T]
    \begin{column}{0.56\textwidth}
\begin{codigo}[title=src/androidTest/.../BancadaDaoTest.kt]{Kotlin}
@RunWith(AndroidJUnit4::class)
class BancadaDaoTest {
  private lateinit var db: AppDatabase
  private lateinit var dao: BancadaDao

  @Before fun criar() {
    db = Room.inMemoryDatabaseBuilder(
      ApplicationProvider.getApplicationContext(),
      AppDatabase::class.java).build()
    dao = db.bancadaDao()
  }
  @After fun fechar() = db.close()

  @Test fun salvarEObservar() = runTest {
    dao.salvar(BancadaEntity(1, "B1", "L-12"))
    val lista = dao.observarTodas().first()
    assertEquals("B1", lista.single().nome)
  }
}
\end{codigo}
    \end{column}
    \begin{column}{0.42\textwidth}
      \begin{conceito}[Banco em memória]
        Criado do zero a cada teste e descartado no fim: testes independentes e rápidos, sem afetar dados reais.
      \end{conceito}
      \begin{dica}[Migrações]
        \cod{MigrationTestHelper} cria o banco na versão antiga, aplica a migração e valida o esquema resultante.
      \end{dica}
      \begin{dica}
        Os testes em \arq{androidTest} rodam com \cod{./gradlew connectedAndroidTest} em um aparelho ou emulador.
      \end{dica}
    \end{column}
  \end{columns}
\end{frame}

\begin{frame}[fragile]{Testes de interface em Compose}
  \begin{columns}[T]
    \begin{column}{0.56\textwidth}
\begin{codigo}[title=TelaPainelTest.kt]{Kotlin}
class TelaPainelTest {
  @get:Rule val compose = createComposeRule()

  @Test fun mostraAlertaQuandoAcimaDoLimite() {
    compose.setContent {
      CartaoBancada(numero = 3, temperatura = 35.0)
    }
    compose.onNodeWithText("Acima do limite!")
      .assertIsDisplayed()
  }

  @Test fun clicarNaBancadaChamaCallback() {
    var aberta = -1
    compose.setContent {
      TelaPainel(estadoExemplo, {}, {}, onAbrir = { aberta = it })
    }
    compose.onNodeWithTag("bancada-2").performClick()
    assertEquals(2, aberta)
  }
}
\end{codigo}
    \end{column}
    \begin{column}{0.42\textwidth}
      \begin{conceito}[Encontrar, agir, verificar]
        \scriptsize
        \textbf{Encontrar}: \cod{onNodeWithText}, \cod{onNodeWithTag}, \cod{onNodeWithContentDescription}.\\
        \textbf{Agir}: \cod{performClick}, \cod{performTextInput}.\\
        \textbf{Verificar}: \cod{assertIsDisplayed}, \cod{assertExists}.
      \end{conceito}
      \begin{conceito}[Telas sem estado]
        A divisão \emph{Route}/\emph{Screen} do Módulo 06 permite testar a tela passando qualquer \cod{UiState}, sem ViewModel nem rede.
      \end{conceito}
    \end{column}
  \end{columns}
\end{frame}

% =============================================================================
\section{Desempenho}

\begin{frame}{Diagnosticando desempenho}
  \begin{columns}[T]
    \begin{column}{0.5\textwidth}
      \begin{conceito}[Ferramentas]
        \begin{itemize}
          \item \textbf{Profiler} do Android Studio: CPU, memória, energia; identifica \emph{jank} (quadros perdidos).
          \item \textbf{Layout Inspector}: contagem de recomposições por composable.
          \item \textbf{Macrobenchmark}: mede inicialização e rolagem de forma reproduzível.
          \item \textbf{Android vitals} (Play Console): ANRs e falhas nos aparelhos reais.
        \end{itemize}
      \end{conceito}
      \begin{atencao}
        Meça sempre em build \textbf{release} e em aparelho real: o build debug é muito mais lento.
      \end{atencao}
    \end{column}
    \begin{column}{0.48\textwidth}
      \begin{dica}[Compose rápido]
        \begin{itemize}
          \item Nada de trabalho pesado durante a composição: formate e filtre no ViewModel.
          \item \cod{key} estável em listas \emph{lazy}.
          \item \cod{derivedStateOf} para estados que mudam com frequência (ex.: posição de rolagem).
          \item Classes de estado imutáveis (\cod{val}, \cod{List}).
        \end{itemize}
      \end{dica}
      \begin{conceito}[Baseline Profiles]
        Arquivo que lista o código usado na inicialização. A Play Store pré-compila esse código e o app abre bem mais rápido.
      \end{conceito}
    \end{column}
  \end{columns}
\end{frame}

% =============================================================================
\section{Acessibilidade}

\begin{frame}[fragile]{Apps para todas as pessoas}
  \begin{columns}[T]
    \begin{column}{0.5\textwidth}
\begin{codigo}[title=Semântica em Compose]{Kotlin}
// imagem informativa: descreva
Icon(Icons.Default.Warning,
  contentDescription = "Temperatura acima do limite")

// card inteiro lido como um único item
Card(Modifier.semantics(mergeDescendants = true) {}) {
  Text("Bancada 3"); Text("24,8 °C")
}

// estado anunciado pelo leitor de tela
Switch(checked = ligado, onCheckedChange = { ligado = it },
  modifier = Modifier.semantics {
    stateDescription = if (ligado) "ligado" else "desligado"
  })
\end{codigo}
    \end{column}
    \begin{column}{0.48\textwidth}
      \begin{conceito}[Checklist]
        \begin{itemize}
          \item \cod{contentDescription} em ícones e imagens informativas.
          \item Alvos de toque de pelo menos \textbf{48\,dp}.
          \item Contraste mínimo de \textbf{4,5:1} para texto.
          \item Funciona com fonte em \textbf{200\,\%} (use \cod{sp}).
          \item Informação não depende só de cor (verde/vermelho + ícone ou texto).
          \item Navegável com \textbf{TalkBack} e por teclado.
        \end{itemize}
      \end{conceito}
      \begin{dica}
        O app \textbf{Accessibility Scanner} analisa a tela aberta e aponta problemas automaticamente.
      \end{dica}
    \end{column}
  \end{columns}
\end{frame}

% =============================================================================
\section{Qualidade contínua}

\begin{frame}[fragile]{Análise estática e integração contínua}
  \begin{columns}[T]
    \begin{column}{0.52\textwidth}
\begin{codigo}[title=.github/workflows/android.yml]{bash}
name: Android CI
on: [push, pull_request]
jobs:
  build:
    runs-on: ubuntu-latest
    steps:
      - uses: actions/checkout@v4
      - uses: actions/setup-java@v4
        with: { distribution: temurin, java-version: 17 }
      - uses: gradle/actions/setup-gradle@v4
      - run: ./gradlew lint testDebugUnitTest assembleDebug
      - uses: actions/upload-artifact@v4
        with:
          name: apk-debug
          path: app/build/outputs/apk/debug/*.apk
\end{codigo}
    \end{column}
    \begin{column}{0.46\textwidth}
      \scriptsize
      \begin{tabular}{@{}l p{3.9cm}@{}}
        \toprule
        \cabecalho{Ferramenta} & \cabecalho{Verifica} \\
        \midrule
        \textbf{Android Lint} & APIs incorretas, permissões, recursos não usados, acessibilidade \\
        \textbf{ktlint} & formatação e estilo Kotlin \\
        \textbf{detekt} & complexidade e más práticas \\
        \textbf{Testes} & comportamento \\
        \bottomrule
      \end{tabular}
      \begin{conceito}[Integração contínua]
        A cada \emph{push}, um servidor compila, analisa e testa o projeto. Um \emph{pull request} só é aceito com tudo verde.
      \end{conceito}
      \begin{dica}
        Cheque as versões atuais das \emph{actions} no GitHub Marketplace antes de usar.
      \end{dica}
    \end{column}
  \end{columns}
\end{frame}

% =============================================================================
\section{Fechamento}

\begin{frame}{Resumo}
  \begin{resumo}
    \begin{itemize}
      \item Muitos testes unitários, alguns de integração e poucos de UI.
      \item \cod{runTest} + \cod{MainDispatcherRule} para coroutines; Turbine para \cod{Flow}.
      \item Fakes para simular dependências; mocks (MockK) para verificar interações.
      \item Room em memória e \cod{createComposeRule} para integração e interface.
      \item Desempenho: meça em release, no aparelho, com Profiler, Layout Inspector e Baseline Profiles.
      \item Acessibilidade e lint fazem parte da qualidade; a CI verifica tudo a cada \emph{push}.
    \end{itemize}
  \end{resumo}
\end{frame}

\begin{frame}{Exercícios}
  \begin{exercicio}[Exercício 1 --- Unitários]
    Escreva testes para \cod{VerificarAlertasUseCase} cobrindo: lista vazia, todas abaixo do limite, uma acima, temperatura nula e mudança do limite.
  \end{exercicio}
  \begin{exercicio}[Exercício 2 --- ViewModel]
    Teste o \cod{PainelViewModel} para os estados carregando, sucesso, vazio e erro de sincronização usando \cod{FakeBancadaRepository} e Turbine.
  \end{exercicio}
  \begin{exercicio}[Exercício 3 --- Acessibilidade]
    Rode o Accessibility Scanner nas telas do LabMonitor, navegue com TalkBack e corrija ao menos três problemas encontrados.
  \end{exercicio}
  \begin{projeto}[LabMonitor --- Etapa 11]
    Atinja ao menos 70\,\% de cobertura nos casos de uso e ViewModels, escreva dois testes de UI para os fluxos de reserva e configure o GitHub Actions para rodar lint e testes.
  \end{projeto}
\end{frame}

\begin{frame}{Para aprofundar}
  \small
  \begin{itemize}
    \item Testes no Android: \url{https://developer.android.com/training/testing}
    \item Testando coroutines: \url{https://developer.android.com/kotlin/coroutines/test}
    \item Testes em Compose: \url{https://developer.android.com/develop/ui/compose/testing}
    \item Turbine: \url{https://github.com/cashapp/turbine} \quad MockK: \url{https://mockk.io}
    \item Desempenho em Compose: \url{https://developer.android.com/develop/ui/compose/performance}
    \item Baseline Profiles: \url{https://developer.android.com/topic/performance/baselineprofiles}
    \item Acessibilidade: \url{https://developer.android.com/guide/topics/ui/accessibility}
  \end{itemize}
\end{frame}

\encerramento{Módulo 12 --- Publicação na Google Play: assinatura, AAB, Play Console, testes fechados, LGPD e CI/CD.}

\end{document}
